\begin{proof}
This proof is similar to the proof
of~\cite[Theorem 3.21]{antichain-toggling}.

Let $(x_1,x_2,\dots,x_n)$ be any linear extension of a finite poset $P$. By the definitions,
$\NAR=\tau_{x_n}\cdots \tau_{x_2}\tau_{x_1}$
and
$\NOR=T_{x_1} T_{x_2} \cdots T_{x_n}$.

Using the isomorphism from $\ntog_A(P)$ to $\ntog_O(P)$ given by $\tau_v \mapsto \tau_v^*$,
it suffices to show that
$\tau_{x_n}^*\cdots \tau_{x_2}^* \tau_{x_1}^*=\NOR = T_{x_1} T_{x_2} \cdots T_{x_n}$.
We will use induction to prove that $\tau_{x_k}^*\cdots \tau_{x_2}^* \tau_{x_1}^*
=T_{x_1} T_{x_2} \cdots T_{x_k}$ for $1\leq k\leq n$.

For the base case, $\tau_{x_1}^*=T_{x_1}$ since $x_1$ is a minimal element of $P$.
For the induction hypothesis, let $1\leq k\leq n-1$ and assume that
$\tau_{x_k}^*\cdots \tau_{x_2}^* \tau_{x_1}^*
=T_{x_1} T_{x_2} \cdots T_{x_k}$.
Then 
\begin{equation}\label{eq:future frenzy}
\tau_{x_{k+1}}^*\tau_{x_k}^*\cdots \tau_{x_2}^* \tau_{x_1}^*
= \eta_{x_{k+1}} T_{x_{k+1}} \eta_{x_{k+1}}^{-1} T_{x_1} T_{x_2} \cdots T_{x_k}.
\end{equation}
Let $(y_1,\dots,y_{k'})$ be a linear extension of the subposet $\{y\in P \;|\; y<x_{k+1}\}$ of $P$.
Then since $(x_1,\dots,x_n)$ is a linear extension of $P$, all of $y_1,\dots,y_{k'}$ must be in $\{x_1,\dots,x_k\}$.
Furthermore, any element less than one of $y_1,\dots,y_{k'}$ must be less than $x_{k+1}$ so none of the elements of $\{x_1,\dots,x_k\}$ outside of $\{y_1,\dots,y_{k'}\}$ are less than any of $y_1,\dots,y_{k'}$.
Therefore, we can name these elements in such a way that $(y_1,\dots,y_{k'}, y_{k'+1}, \dots, y_k)$ is a linear extension of $\{x_1,\dots,x_k\}$.
Noting again that any two linear extensions of a poset differ by a sequence of swaps between
adjacent incomparable elements~\cite{etienne-84}, 
toggles of incomparable elements commute so $T_{x_1}T_{x_2}\cdots T_{x_k} = T_{y_1}\cdots T_{y_{k'}}T_{y_{k'+1}}\cdots T_{y_k}$.
From Eq.~\eqref{eq:future frenzy} and $\eta_{x_{k+1}}=T_{y_1}\cdots T_{y_{k'}}$, we obtain
\begin{align*}
\tau_{x_{k+1}}^*\tau_{x_k}^*\cdots \tau_{x_2}^* \tau_{x_1}^* &=
\eta_{x_{k+1}} T_{x_{k+1}} \eta_{x_{k+1}}^{-1}
T_{x_1} T_{x_2} \cdots T_{x_k}
\\&=
T_{y_1}\cdots T_{y_{k'}} T_{x_{k+1}} E_{y_{k'}}\cdots E_{y_1} T_{y_1}\cdots T_{y_{k'}}T_{y_{k'+1}}\cdots T_{y_k}
\\&=
T_{y_1}\cdots T_{y_{k'}} T_{x_{k+1}} T_{y_{k'+1}}\cdots T_{y_k}
\\&=
T_{y_1}\cdots T_{y_{k'}} T_{y_{k'+1}}\cdots T_{y_k} T_{x_{k+1}}
\\&=
T_{x_1} T_{x_2} \cdots T_{x_k} T_{x_{k+1}}.
\end{align*}
In the fourth equality above, we could move $T_{x_{k+1}}$ to the right of $T_{y_{k'+1}}\cdots T_{y_k}$ because $x_{k+1}$ is incomparable with each of $y_{k'+1},\dots, y_k$. This is because none of these are less than $x_{k+1}$ by design nor greater than $x_{k+1}$ by position within the linear extension $(x_1,\dots,x_n)$ of $P$.

By induction, we have $\tau_{x_n}^*\cdots \tau_{x_2}^* \tau_{x_1}^* = T_{x_1} T_{x_2} \cdots T_{x_n} = \NOR = \Theta\circ\up^{-1}\circ \down$ so
\[
\tau_{x_n}\cdots \tau_{x_2} \tau_{x_1} = \NAR
=\down\circ\Theta\circ\up^{-1}.\qedhere
\]
\end{proof}
